%%
%% Copyright 2019-2024 Elsevier Ltd
%%
%% This file is part of the 'CAS Bundle'.
%% --------------------------------------
%%
%% Manuscript for Energy Conversion & Management
%% Using cas-dc (double column) template
%%

\documentclass[a4paper,fleqn]{cas-dc}

\usepackage[numbers]{natbib}
\usepackage{amsmath,amssymb}
\usepackage{booktabs}
\usepackage{graphicx}
\usepackage{siunitx}
\usepackage{multirow}
\usepackage{lineno}

\sisetup{
  per-mode=symbol,
  inter-unit-product=\ensuremath{\cdot},
}

%%% Author macros
\def\tsc#1{\csdef{#1}{\textsc{\lowercase{#1}}\xspace}}
\tsc{WGM}
\tsc{QE}
%%%

\newtheorem{theorem}{Theorem}
\newtheorem{lemma}[theorem]{Lemma}

\begin{document}
\let\WriteBookmarks\relax
\def\floatpagepagefraction{1}
\def\textpagefraction{.001}

% Short title
\shorttitle{Network Topology Abstraction in Heat Network Dispatch Optimization}

% Short author
\shortauthors{L.\ Ruess et~al.}

% Main title
\title[mode=title]{Effect of Network Topology Abstraction on Operational Dispatch Optimization of Electrified Industrial Heat Networks: A MILP Comparative Study}

% First author
\author[1]{Lukas Ruess}
\cormark[1]
\ead{lukas.ruess@fhv.at}
\credit{Conceptualization, Methodology, Software, Validation, Formal analysis, Writing -- original draft, Visualization}

\affiliation[1]{organization={Research Centre Energy, Fachhochschule Vorarlberg},
            addressline={Hochschulstra{\ss}e 1},
            city={Dornbirn},
            postcode={6850},
            state={Vorarlberg},
            country={Austria}}

% Second author — TODO: fill in
\author[1]{Second Author}
\credit{Supervision, Writing -- review \& editing}

% Third author — TODO: fill in
\author[1]{Third Author}
\credit{Resources, Writing -- review \& editing}

% Corresponding author text
\cortext[1]{Corresponding author}

% Abstract
\begin{abstract}
District heating system planners must choose how much spatial detail to include in network optimization models, yet the literature provides little quantitative guidance on this trade-off because prior comparisons simultaneously change both the loss physics and the topology.
This paper fills that gap by comparing three levels of network topology abstraction---copperplate (L1, single node), simplified 5-node (L2), and realistic 29-node (L3)---using \emph{identical} physics-based thermal loss models across all three levels, so that topology is the sole experimental variable.
The comparison is applied to the Stadtbach industrial heat network in Austria (\SI{76}{MW} peak demand, \SI{517}{GWh/year}) using a mixed-integer linear programming (MILP) formulation that incorporates piecewise-linear pipe heat losses, pre-computed temperature-dependent heat pump coefficient of performance (COP), and stratified storage geometry.
Results show that L1's zero-loss copperplate assumption underestimates annual operational cost by 0.5\% (\EUR{74,195}) compared to L3, while L2 captures over 85\% of the L1--L3 cost difference at approximately 40\% shorter solve time.
For this network with homogeneous industrial demand, topology detail beyond L2 provides marginal operational value.
The optimizer selects a fully electrified dispatch (zero gas boiler usage) across all scenarios at a mean heat pump COP of 4.0.
Practitioner guidance is provided: L1 suffices for screening studies, L2 for detailed planning, and L3 for final engineering validation.
\end{abstract}

% Research highlights
\begin{highlights}
\item First controlled three-level topology comparison (L1/L2/L3) with identical loss physics in MILP dispatch optimization
\item Physics-based piecewise-linear pipe loss model integrated into tractable MILP formulation
\item Simplified 5-node model captures 85\% of full 29-node model cost correction at 40\% shorter solve time
\item COP pre-computation linearization introduces $\leq$1.4\% total cost error with formal error bound
\item Full-year 8,760-hour MILP solved in under 20 minutes on commodity hardware
\end{highlights}

% Keywords
\begin{keywords}
Mixed-integer linear programming \sep
District heating optimization \sep
Network topology abstraction \sep
Piecewise-linear thermal loss model \sep
Operational dispatch optimization \sep
Heat pump COP linearization
\end{keywords}

\maketitle

% Line numbering for review
\linenumbers


%% ====================================================================
\section{Introduction}\label{sec:introduction}
%% ====================================================================

\subsection{Motivation and problem context}

Industrial heat supply has emerged as a critical bottleneck in Europe's decarbonization pathway. The heating sector accounts for approximately 50\% of final energy consumption across the EU-27, with industry and district heating together representing 28\% of total demand~\citep{EC2021,IEA2022}. Traditional fossil-fuel-based heat production---particularly natural gas boilers and combined heat and power (CHP) plants---remains the dominant technology, responsible for roughly 60\% of industrial thermal energy generation and 35--40\% of energy-related CO\textsubscript{2} emissions from manufacturing~\citep{EEA2023}.

Transitioning to electrified heating networks requires district heating systems to integrate multiple heat sources---heat pumps, waste heat recovery, thermal storage---across spatially distributed pipe networks~\citep{Werner2017,Paardekooper2018}. When planning or operating such systems, practitioners must choose how much spatial detail to include in their network model. This seemingly technical choice has direct consequences for optimization accuracy, computational cost, and ultimately for the quality of dispatch decisions.

A critical but under-quantified question arises at every planning study: \emph{Does network topology detail change operational outcomes and cost when the underlying loss physics are held constant?}

Practitioners face a common dilemma. Simplified approaches (copperplate or aggregated nodes) reduce computational burden but may miss spatial constraints and zonal demand imbalances. Detailed approaches (20--30 nodes with explicit pipe routing) capture realistic spatial structure but increase model size and solve time. Surprisingly, the literature provides little quantitative guidance on this trade-off, because most prior comparisons also change the loss physics model between topology levels---making it impossible to isolate the topology effect from the loss modeling effect~\citep{vanderHeijde2017,Li2013}.

This paper fills that gap. We compare three levels of network topology abstraction---copperplate (Level~1, L1), simplified 5-node (Level~2, L2), and detailed 29-node (Level~3, L3)---using identical physics-based thermal loss models across all three, so that network topology is the sole experimental variable. We apply these three models to a real industrial heating network in Austria using one year of operational data, and we quantify: (i)~operational cost differences attributable purely to topology abstraction; (ii)~dispatch pattern differences under different network structures; and (iii)~computational cost as a function of model resolution.

\subsection{Network topology classification}

The literature reveals a spectrum of thermal network modeling approaches. This work classifies network abstraction into three levels:

\textbf{L1: Copperplate (aggregated single node).} All components and demand at a single virtual node with no spatial differentiation. No pipe losses modeled. Fast computation (2--3\,min for full-year optimization) but ignores network losses entirely, potentially underestimating system cost by 5--10\%~\citep{Leitner2018}.

\textbf{L2: Simplified multi-node (aggregated zones).} Network represented as 5--10 strategically placed nodes with aggregated demand zones. Uses identical physics-based loss model to L3 ($Q_\mathrm{loss} = U \cdot L \cdot \Delta T$ with piecewise linearization) but simplified topology. Captures spatial demand distribution at moderate computational cost.

\textbf{L3: Realistic multi-node (full network detail).} Actual network topology with 20--30 nodes and realistic pipe routing. Same physics-based loss model as L2. Captures local constraints and flow patterns suitable for detailed engineering design.

Unlike prior work, this study uses identical thermal physics (piecewise-linear loss model) across all three levels, isolating network topology from loss modeling fidelity.

\subsection{Research gap and contributions}

The literature provides extensive treatments of MILP-based energy system optimization~\citep{Lund2009,Bloess2018,Robinius2017} and of district heating network simulation~\citep{Benonysson1995,Frederiksen2013,Moller2017}, but a methodological gap persists: no prior study quantifies the impact of network topology abstraction in isolation, holding all other modeling choices constant. Existing comparisons simultaneously change the loss physics, the number of components, and sometimes the solver~\citep{Moser2019,Wang2020,Volkova2019,Leitner2018}.

Furthermore, the open-source frameworks most commonly used for dispatch optimization---oemof~\citep{Hilpert2018}, PyPSA~\citep{Brown2018}, TIMES~\citep{Loulou2005}---treat thermal networks as single-node copperplate systems.

This paper addresses this gap through five contributions:
\begin{enumerate}
\item \textbf{Topology abstraction study} (primary): A controlled three-scenario comparison (L1/L2/L3) on a real industrial heating network in which identical piecewise-linear pipe loss models are applied at all levels.
\item \textbf{MILP formulation}: A complete mixed-integer linear program for operational dispatch incorporating physical pipe heat losses, temperature-dependent COP via pre-computed time series, and piecewise-linear stratified storage losses.
\item \textbf{Linearization with error bounds}: Pre-computing COP as an hourly time series enables MILP dispatch without bilinear terms, with a formal error bound of $\pm$2--3\% on system cost.
\item \textbf{Open-source implementation}: The CALION framework (Python/Pyomo) implements all three topology levels with sensitivity analysis and benchmarking modules.
\item \textbf{Practitioner guidance}: Quantitative evidence for when L1/L2/L3 resolution is sufficient for different planning contexts.
\end{enumerate}

\subsection{Paper organization}

Section~\ref{sec:literature} reviews relevant literature. Section~\ref{sec:methodology} develops the mathematical formulation. Section~\ref{sec:casestudy} describes the Stadtbach case study. Section~\ref{sec:results} presents results. Section~\ref{sec:discussion} discusses implications and limitations. Section~\ref{sec:conclusion} concludes.


%% ====================================================================
\section{Literature review}\label{sec:literature}
%% ====================================================================

\subsection{MILP-based energy system optimization}

Mixed-integer linear programming has become the standard paradigm for capacity and dispatch optimization in energy systems. Lund and Mathiesen~\citep{Lund2009} used MILP to study the role of district heating and storage in 100\% renewable grids. Multi-energy system MILP frameworks have advanced considerably: Mancarella~\citep{Mancarella2014} established multi-energy hub modeling; Gabrielli et al.~\citep{Gabrielli2018} formulated MILP for distributed systems including thermal and electrical storage.

Open-source tools have democratized energy system optimization. \textbf{oemof}~\citep{Hilpert2018} provides a Python-based MILP core but treats district heating with simplified aggregate efficiency factors. \textbf{PyPSA}~\citep{Brown2018} offers MILP solvers with multi-carrier support but limited thermal spatial detail. \textbf{Calliope}~\citep{Pfenninger2018} supports multi-node modeling but represents pipe physics via aggregate efficiencies. None of these frameworks model physical pipe heat loss at the individual pipe level within the optimization.

\subsection{District heating network modeling}

Benonysson et al.~\citep{Benonysson1995} pioneered combined hydraulic-thermal models for district heating. The standard pipe heat loss formula $Q_\mathrm{loss} = U \cdot L \cdot \Delta T_\mathrm{mean}$ is well-validated empirically~\citep{Frederiksen2013}. The overall heat transfer coefficient $U$ for pre-insulated pipes typically ranges from 0.15--0.40\,\si{W/(m.K)}~\citep{vanderHeijde2017}.

Several recent works have incorporated multi-node representations into planning studies. Moser et al.~\citep{Moser2019} developed a graph-based MILP for district heating expansion but using constant efficiency factors. Wang et al.~\citep{Wang2020} employed a nodal temperature model with heuristic dispatch. Volkova et al.~\citep{Volkova2019} found that pipe routing significantly affects loss estimates, but in a simulation context. Leitner et al.~\citep{Leitner2018} noted that copperplate assumptions underestimate costs by 3--8\%, though without a controlled topology comparison.

None of these studies isolates network topology as the sole experimental variable.

\subsection{Linearization methods for MILP integration}

Piecewise-linear (PWL) approximations of smooth functions~\citep{Rebennack2016} enable integration of physical nonlinearities into MILP. For COP pre-computation, Wirtz et al.~\citep{Wirtz2018} demonstrated that pre-computing heat pump COP from source temperature data enables accurate MILP dispatch. Ommen et al.~\citep{Ommen2016} showed that COP sensitivity to part-load ratio ($\pm$5\%) is secondary compared to source temperature effects ($\pm$30\%), supporting pre-computation at nominal load.

\subsection{Positioning of this work}

Table~\ref{tab:positioning} summarizes the gap addressed. The MILP optimization community (oemof, PyPSA, TIMES) treats thermal networks as copperplate (L1). The district heating simulation community (TRNSYS, Modelica) offers physical accuracy but no optimization capability. Our study quantifies how much the copperplate assumption costs by providing the first controlled comparison with identical physics.

\begin{table}[t]
\centering
\caption{Positioning relative to related work. The present study is the first to isolate topology as the sole variable across three abstraction levels with MILP optimization.}\label{tab:positioning}
\begin{tabular*}{\tblwidth}{@{}lcccc@{}}
\toprule
Study & Levels & Same physics? & MILP? & Gap addressed \\
\midrule
Moser et al.~\citep{Moser2019} & 1 & --- & \checkmark & No L1/L2 comparison \\
Wang et al.~\citep{Wang2020} & 1 & --- & $\times$ & No MILP \\
Volkova et al.~\citep{Volkova2019} & 2 & $\times$ & $\times$ & Physics varies \\
Leitner et al.~\citep{Leitner2018} & 2 & $\times$ & \checkmark & Physics varies \\
\textbf{This study} & \textbf{3} & \textbf{\checkmark} & \textbf{\checkmark} & \textbf{Topology isolated} \\
\bottomrule
\end{tabular*}
\end{table}


%% ====================================================================
\section{Methodology}\label{sec:methodology}
%% ====================================================================

\subsection{System description and scope}

We consider an industrial heat network comprising: (i)~heat generation from multiple technologies (heat pumps, CHP, gas boilers); (ii)~a stratified thermal storage tank; (iii)~a pipe distribution network with distance-dependent heat losses; (iv)~industrial process heat demand; and (v)~external coupling to the electricity grid and waste heat sources.

The system boundary covers one calendar year (8,760 hours at hourly resolution). The case study in Section~\ref{sec:casestudy} holds all asset capacities fixed (dispatch-only mode) so that network topology is the sole experimental variable.

\subsection{Mathematical formulation}

\subsubsection{Sets, parameters, and decision variables}

The temporal set is $\mathcal{T} = \{1, 2, \ldots, 8760\}$ with $\Delta t = 1$\,h. Component sets include generators $\mathcal{G}$ (with CHP subset $\mathcal{G}_\mathrm{CHP} \subset \mathcal{G}$) and heat pumps. Economic parameters include time-varying electricity price $p_\mathrm{el}[t]$\,[\EUR{}/MWh], grid emissions factor $\mathrm{ef}_\mathrm{grid}[t]$\,[kg\,CO\textsubscript{2}/MWh], and fuel prices $p_f(g)$\,[\EUR{}/MWh]. Technical parameters include efficiencies $\eta_{\mathrm{th},g}$, $\eta_{\mathrm{el},g}$, pre-computed COP$[t]$, storage loss rate $\lambda_\mathrm{loss}$, and charge/discharge efficiencies $\eta_c$, $\eta_d$.

Decision variables comprise: continuous hourly variables for generation ($Q_g[t]$, $F_g[t]$, $P_{\mathrm{el},g}[t]$), heat pump ($Q_\mathrm{hp}[t]$, $P_\mathrm{hp}[t]$), grid ($P_\mathrm{buy}[t]$, $P_\mathrm{sell}[t]$), storage ($E_\mathrm{tes}[t]$, $Q_c[t]$, $Q_d[t]$), and flexibility ($Q_\mathrm{dump}[t]$); binary variables for control ($y_{\mathrm{on},g}[t]$, $y_\mathrm{charge}[t]$, $y_\mathrm{buy}[t]$).

Total model size: approximately 44,000--56,000 continuous and 26,000--30,000 binary variables for the 1-year horizon, depending on topology level.

\subsubsection{Energy balance constraints}

The heat balance requires that supply equals demand plus losses at each time step:
\begin{equation}
\sum_{g \in \mathcal{G}} Q_g[t] + Q_\mathrm{hp}[t] + Q_d[t] = Q_\mathrm{dem}[t] + Q_c[t] + Q_\mathrm{loss}[t] + Q_\mathrm{dump}[t] \quad \forall t
\label{eq:heat_balance}
\end{equation}

The electricity balance couples generation and consumption:
\begin{equation}
P_\mathrm{buy}[t] + \sum_{g \in \mathcal{G}_\mathrm{CHP}} P_{\mathrm{el},g}[t] = P_\mathrm{sell}[t] + P_\mathrm{hp}[t] \quad \forall t
\label{eq:elec_balance}
\end{equation}

\subsubsection{Thermal generation constraints}

For each generator $g$:
\begin{align}
Q_g[t] &= \eta_{\mathrm{th},g} \cdot F_g[t] \label{eq:gen_eff} \\
P_{\mathrm{el},g}[t] &= \eta_{\mathrm{el},g} \cdot F_g[t] \quad (g \in \mathcal{G}_\mathrm{CHP}) \label{eq:chp_eff} \\
\lambda_\mathrm{min} \cdot \mathrm{cap}_g \cdot y_{\mathrm{on},g}[t] &\leq Q_g[t] \leq \mathrm{cap}_g \cdot y_{\mathrm{on},g}[t] \label{eq:gen_bounds}
\end{align}

\subsubsection{Heat pump constraints (Linearization \#1)}

The COP-based output constraint is:
\begin{equation}
Q_\mathrm{hp}[t] = \mathrm{COP}[t] \cdot P_\mathrm{hp}[t] \quad \forall t
\label{eq:hp_cop}
\end{equation}
where $\mathrm{COP}[t]$ is a pre-computed scalar parameter (not a decision variable). This formulation is linear in $P_\mathrm{hp}[t]$, making it compatible with MILP solvers (see Section~\ref{sec:linearization_cop}).

\paragraph{Minimum uptime and downtime.}
To prevent unrealistic cycling of the compressor, startup $u[t]\in\{0,1\}$ and shutdown $v[t]\in\{0,1\}$ auxiliary binaries are introduced following~\citet{Carrion2006}:
\begin{align}
y[t] - y[t{-}1] &= u[t] - v[t] \label{eq:hp_transition} \\
u[t] + v[t]     &\leq 1 \label{eq:hp_no_simult} \\
\sum_{k=t-L+1}^{t} u[k] &\leq y[t] \quad \forall t \label{eq:hp_uptime} \\
\sum_{k=t-D+1}^{t} v[k] &\leq 1 - y[t] \quad \forall t \label{eq:hp_downtime}
\end{align}
where $L$~[timesteps] is the minimum uptime and $D$~[timesteps] the minimum downtime. Setting $L = D = 0$ recovers the unconstrained binary on/off model.

\subsubsection{Thermal storage constraints (Linearization \#2)}

State of charge dynamics:
\begin{equation}
E_\mathrm{tes}[t] = E_\mathrm{tes}[t{-}1](1 - \lambda_\mathrm{loss}) + \eta_c Q_c[t] \Delta t - \frac{Q_d[t]}{\eta_d} \Delta t - Q_\mathrm{loss,tes}[t]
\label{eq:storage_soc}
\end{equation}
where $Q_\mathrm{loss,tes}[t]$ is computed via piecewise-linear approximation (Section~\ref{sec:linearization_storage}). Mutual exclusivity of charge/discharge is enforced via Big-M:
\begin{align}
Q_c[t] &\leq P_\mathrm{max} \cdot y_\mathrm{charge}[t] \label{eq:stor_charge} \\
Q_d[t] &\leq P_\mathrm{max} \cdot (1 - y_\mathrm{charge}[t]) \label{eq:stor_discharge}
\end{align}

\subsubsection{Network loss constraints}

Physical heat loss for each pipe $i$:
\begin{equation}
Q_{\mathrm{loss},i}[t] = \frac{U_i \cdot L_i}{1000} \cdot (T_s - T_\mathrm{amb}) \quad [\mathrm{MW}]
\label{eq:pipe_loss}
\end{equation}
where $U_i$\,[\si{W/(m.K)}] is the heat transfer coefficient, $L_i$\,[m] the pipe length, $T_s$\,[\si{\degreeCelsius}] the supply temperature, and $T_\mathrm{amb}$\,[\si{\degreeCelsius}] the ground temperature. For L1, $Q_\mathrm{loss}[t] = 0$. For L2 and L3, losses are summed over all pipes in the respective topology.

\paragraph{Pump power model.}
Each pipe segment is equipped with a circulation pump. Pump electricity consumption is approximated as a piecewise-linear (PWL) function of mass flow, using the same three-segment breakpoints as the Darcy--Weisbach pressure drop:
\begin{equation}
P_{\mathrm{pump},i}[t] \approx \sum_{k=1}^{3}
  \bigl(s_k^P \cdot \dot{m}_{i,k}[t] + b_k^P \cdot z_{i,k}[t]\bigr) \quad [\text{MW}]
\label{eq:pump_power}
\end{equation}
where slopes $s_k^P$ and intercepts $b_k^P$ are pre-computed from:
\begin{equation}
P_{\mathrm{pump},i,\ell} =
  \frac{2\,\Delta p_{i,\ell} \cdot \dot{m}_{i,\ell} \cdot 10^5}
       {\rho \cdot \eta_{\mathrm{pump}} \cdot 10^6}
\quad [\text{MW}], \quad \eta_{\mathrm{pump}} = 0.70
\label{eq:pump_bp}
\end{equation}
with $\Delta p_{i,\ell}$ the supply pressure drop at breakpoint $\ell$ and $\rho = 1000$\,\si{kg/m^3}. The nodal pressure balance at each producer node enforces:
\begin{equation}
p_j^{\mathrm{ret}}[t] = p_j^{\mathrm{sup}}[t]
  - \Delta p_i^{\mathrm{sup}}[t] - \Delta p_i^{\mathrm{ret}}[t]
\label{eq:pump_head}
\end{equation}
Pump power $\sum_i P_{\mathrm{pump},i}[t]$ is added to the electricity demand bus, increasing the accuracy of electricity cost estimates.

\subsubsection{Physical state constraints}

In multi-node models (L2, L3), physical realizability is enforced:
\begin{align}
T_{\mathrm{supply},n}[t] &\geq T_{\mathrm{return},n}[t] - \epsilon_T \quad \forall n \in \mathcal{N} \label{eq:temp_valid} \\
p_n[t] &\geq p_\mathrm{min} \quad \forall n \in \mathcal{N} \label{eq:pressure_min} \\
v_i[t] &\leq v_\mathrm{max} \quad \forall i \in \mathcal{P} \label{eq:velocity_max}
\end{align}
with $\epsilon_T = \SI{0.1}{\degreeCelsius}$, $p_\mathrm{min} = 0.5$\,bar, and $v_\mathrm{max} = \SI{2.5}{m/s}$. All constraints are linear and preserve MILP tractability.

\subsubsection{Grid coupling}

Mutual exclusivity of grid import/export:
\begin{align}
P_\mathrm{buy}[t] &\leq M \cdot y_\mathrm{buy}[t] \label{eq:grid_buy} \\
P_\mathrm{sell}[t] &\leq M \cdot (1 - y_\mathrm{buy}[t]) \label{eq:grid_sell}
\end{align}
where $M = 10{,}000$\,MW. Peak import tracking for the demand charge: $P_\mathrm{grid,max} \geq P_\mathrm{buy}[t] \;\; \forall t$.

\subsubsection{Objective function}

Minimize total annualized operational cost:
\begin{equation}
Z = C_\mathrm{fuel} + C_\mathrm{elec} + C_{\mathrm{CO}_2} + C_\mathrm{dump} + C_\mathrm{demand}
\label{eq:objective}
\end{equation}
with:
\begin{align}
C_\mathrm{fuel} &= \textstyle\sum_{t} \sum_{g} p_f(g) \cdot F_g[t] \label{eq:cost_fuel} \\
C_\mathrm{elec} &= \textstyle\sum_{t} \bigl[(p_\mathrm{el}[t] + c_\mathrm{fee}) P_\mathrm{buy}[t] - (p_\mathrm{el}[t] - c_\mathrm{spread}) P_\mathrm{sell}[t]\bigr] \label{eq:cost_elec} \\
C_{\mathrm{CO}_2} &= \textstyle\frac{p_{\mathrm{CO}_2}}{1000} \sum_{t} \bigl[\sum_{g} \mathrm{ef}_f(g) F_g[t] + \mathrm{ef}_\mathrm{grid}[t] P_\mathrm{buy}[t]\bigr] \label{eq:cost_co2}
\end{align}

\subsection{Linearization Strategy \#1: COP pre-computation}\label{sec:linearization_cop}

The original nonlinear coupling $Q_\mathrm{hp}[t] = \mathrm{COP}(T_\mathrm{source}[t], T_\mathrm{sink}[t]) \cdot P_\mathrm{hp}[t]$ is resolved by computing COP offline as a function of \emph{both} source and sink temperature. The sink temperature follows the heating curve:
\begin{equation}
T_{\mathrm{sup,hc}}[t] = T_{\mathrm{sup,min}} +
  (T_{\mathrm{sup,max}} - T_{\mathrm{sup,min}}) \cdot
  \frac{T_{\mathrm{amb,high}} - T_{\mathrm{amb}}[t]}
       {T_{\mathrm{amb,high}} - T_{\mathrm{amb,low}}}
\label{eq:heating_curve}
\end{equation}
clamped to $[T_{\mathrm{sup,min}}, T_{\mathrm{sup,max}}]$. Two COP methods are supported:

\textbf{Method~A (analytical):} For a heat pump following the Carnot cycle with practical efficiency $\eta_\mathrm{Carnot} \in [0.3, 0.7]$:
\begin{equation}
\mathrm{COP}[t] = \eta_\mathrm{Carnot} \cdot \frac{T_{\mathrm{sup,hc}}[t]\,[\mathrm{K}]}{T_{\mathrm{sup,hc}}[t]\,[\mathrm{K}] - T_\mathrm{source}[t]\,[\mathrm{K}]}
\label{eq:cop_analytical}
\end{equation}

\textbf{Method~B (tabular):} Bilinear interpolation of a 2D manufacturer COP table at $(T_\mathrm{source}[t],\, T_{\mathrm{sup,hc}}[t])$.

Once $\mathrm{COP}[t]$ becomes a fixed parameter series, Eq.~\eqref{eq:hp_cop} is linear in $P_\mathrm{hp}[t]$, preserving MILP tractability. The heating-curve correction captures COP improvement at lower supply temperatures while avoiding bilinear terms.

\textbf{Error bound.} Interpolation error $\leq 3.6$\% (5-point grid, 30\,K range). Combined with manufacturer table uncertainty ($\pm$2\%, ISO~13256), the total COP error is $\varepsilon_\mathrm{COP} \approx 4.1$\%. For heat pump costs at 25--40\% of OPEX, this propagates to $\leq$1.4\% total system cost error---acceptable for planning studies.

\subsection{Linearization Strategy \#2: PWL storage losses}\label{sec:linearization_storage}

For a stratified cylindrical tank, heat loss depends nonlinearly on fill level through the geometry-dependent surface area $A(h)$. We discretize the loss curve into $N$ piecewise-linear segments over $[0, E_\mathrm{max}]$.

Error bound (standard PWL theory~\citep{Rebennack2016}):
\begin{equation}
\|Q_\mathrm{loss} - Q_\mathrm{loss,PWL}\|_\infty \leq \frac{M (E_\mathrm{max}/N)^2}{8}
\label{eq:pwl_error}
\end{equation}
where $M = \max |Q''_\mathrm{loss}|$. For a \SI{500}{MWh} tank with $N = 10$ segments: error $\leq \SI{31}{W} \approx 0.3$\% of peak loss.

\subsection{Model assumptions and simplifications}\label{sec:assumptions}

Table~\ref{tab:assumptions} summarises accepted simplifications and their justification.

\begin{table}[h]
\centering
\caption{Accepted modelling simplifications and their justification.}
\label{tab:assumptions}
\begin{tabular}{p{4.2cm}p{5.0cm}p{2.8cm}}
\toprule
Assumption & Justification & Max.\ error \\
\midrule
Arithmetic mean for pipe heat loss
  ($\bar{T} = (T_\mathrm{in}+T_\mathrm{out})/2$)
  & MILP compatibility; LMTD introduces $\ln(\cdot)$; standard in DH literature~\citep{Frederiksen2013}
  & ${<}2\%$ for $L{<}2$\,km, $\Delta T{<}5$\,K \\
Constant fluid properties ($c_p$, $\rho$)
  & Variation ${<}6\%$ over 40--120\,\si{\degreeCelsius}
  & ${<}6\%$ mass flow \\
Pre-computed $\mathrm{COP}(t)$ parameter
  & MILP compatibility; corrected via heating curve $T_{\mathrm{sup,hc}}(t)$ (Eq.~\ref{eq:heating_curve})
  & Residual ${<}5\%$ \\
Radial network topology
  & Greenfield planning scope; ring nets: future work
  & N/A \\
Perfect foresight
  & Annual planning standard~\citep{Frederiksen2013}
  & N/A \\
\bottomrule
\end{tabular}
\end{table}

\subsection{Framework classification}

The CALION formulation occupies a position between energy-only MILP tools and full thermo-hydraulic simulation. It includes temperature-dependent COP (via heating-curve sink temperature, Eq.~\ref{eq:heating_curve}), PWL storage geometry, physical pipe heat loss per Eq.~\eqref{eq:pipe_loss}, pump power consumption per Eq.~\eqref{eq:pump_power}, physical state constraints (Eqs.~\ref{eq:temp_valid}--\ref{eq:velocity_max}), and maintains MILP tractability (15--20\,min for L3 over 8,760 hours). It intentionally excludes dynamic mass-flow in pipes and sub-hourly transients (Table~\ref{tab:assumptions}).


%% ====================================================================
\section{Case study: Stadtbach industrial heat network}\label{sec:casestudy}
%% ====================================================================

\subsection{System description}

The case study evaluates the Stadtbach industrial thermal network in the Vorarlberg region of Austria---an existing system serving multiple industrial facilities. The study analyzes operational dispatch on the existing fixed asset base to isolate the network topology effect.

System characteristics: peak demand \SI{76}{MW}; annual heat demand \SI{517}{GWh}; average load factor 77\%. The installed asset base comprises a \SI{200}{MW} gas boiler ($\eta_\mathrm{th} = 0.90$), a \SI{20}{MW} CHP unit (electrical efficiency 18\%, thermal 72.5\%), a \SI{100}{MW} heat pump (air-source + waste heat recovery), a \SI{500}{MWh}/\SI{50}{MW} thermal storage tank, and a \SI{200}{MW} grid connection.

Network infrastructure: total pipe length $\sim$\SI{14250}{m} (pre-insulated dual pipes), pipe $U$-value \SI{0.15}{W/(m.K)}, supply temperature \SI{90}{\degreeCelsius}, return \SI{55}{\degreeCelsius}, ground \SI{10}{\degreeCelsius}.

\subsection{Topology scenarios}

Three topology configurations are compared (Table~\ref{tab:topology}). All three optimize dispatch for identical installed capacities (no investment decisions).

\begin{table}[t]
\centering
\caption{Network topology scenario configuration.}\label{tab:topology}
\begin{tabular*}{\tblwidth}{@{}lccc@{}}
\toprule
Aspect & \textbf{L1} & \textbf{L2} & \textbf{L3} \\
\midrule
Network nodes & 1 & 5 & 29 \\
Consumer nodes & 0 & 3 & 23 \\
Junction nodes & 0 & 1 & 5 \\
Pipes & 0 & 4 & 28 \\
Total pipe length [m] & 0 & 6,300 & 14,250 \\
Pipe diameter [mm] & --- & 350--500 & 150--600 \\
Loss model & $Q_\mathrm{loss}=0$ & PWL (10 seg.) & PWL (10 seg.) \\
Demand distribution & Aggregated & 4 zones & 23 zones \\
\bottomrule
\end{tabular*}
\end{table}

\textbf{L1 (Copperplate):} Single virtual node with all assets and aggregated demand. No pipes, no losses. Provides theoretical lower bound on cost.

\textbf{L2 (Simplified 5-node):} Central plant hub + 4 demand zones connected by 4 trunk lines totaling \SI{6300}{m}. Same PWL loss model as L3.

\textbf{L3 (Realistic 29-node):} 1 plant + 5 junction nodes + 23 consumer zones connected by 28 pipes totaling \SI{14250}{m}. Pipe diameters range from 150--\SI{600}{mm}.

\subsection{Heat pump COP configuration}

All three scenarios use the identical analytical COP method (Eq.~\ref{eq:cop_analytical}) with $T_\mathrm{sink} = \SI{363}{K}$ (\SI{90}{\degreeCelsius}), $T_\mathrm{source}[t]$ from hourly waste heat data (15--\SI{50}{\degreeCelsius}), and $\eta_\mathrm{rel} = 0.6$. The resulting mean COP is approximately 4.0. The identical COP time series across L1/L2/L3 isolates network topology as the only variable.

\subsection{Optimization configuration}

Common settings for all scenarios: dispatch-only mode (no capacity decisions), full-year 2023 horizon (8,760 hours), HiGHS solver~\citep{HiGHS2021} (open-source MIP), 1\% MIP gap, \SI{3600}{s} time limit.


%% ====================================================================
\section{Results}\label{sec:results}
%% ====================================================================

\subsection{Operational cost comparison}

Table~\ref{tab:cost_breakdown} presents the annual cost breakdown by scenario.

\begin{table}[t]
\centering
\caption{Annual cost breakdown by topology scenario (dispatch-only, fixed assets). All values in k\EUR{}.}\label{tab:cost_breakdown}
\begin{tabular*}{\tblwidth}{@{}lrrrl@{}}
\toprule
Category & \textbf{L1} & \textbf{L2} & \textbf{L3} & $\Delta$L3--L1 \\
\midrule
Grid electricity & 9,462 & 9,505 & 9,512 & +0.5\% \\
CO\textsubscript{2} cost & 5,311 & 5,331 & 5,335 & +0.5\% \\
Gas fuel & 0 & 0 & 0 & --- \\
\midrule
\textbf{Total cost} & \textbf{14,772} & \textbf{14,836} & \textbf{14,847} & \textbf{+0.5\%} \\
\bottomrule
\end{tabular*}
\end{table}

Key findings:

\begin{enumerate}
\item \textbf{Gas boiler not dispatched.} The optimizer assigns zero gas consumption across all scenarios. The heat pump (COP~$\approx$~4.0) is strictly more cost-effective than the gas boiler for all hours, leading to a fully electrified dispatch.

\item \textbf{Total cost difference is modest.} L3 costs \EUR{74,195} more annually than L1 (+0.5\%), driven entirely by the additional electricity needed to compensate for pipe heat losses via the heat pump.

\item \textbf{L2 captures most of the L3 effect.} L2's cost is \EUR{63,516} above L1, representing 86\% of the L1--L3 gap (\EUR{74,195}). The additional spatial resolution from L2 to L3 contributes only \EUR{10,679} (0.07\% of total cost).

\item \textbf{CO\textsubscript{2} emissions follow cost.} Total emissions range from 53,108\,t (L1) to 53,349\,t (L3), reflecting the additional grid electricity needed for loss compensation.
\end{enumerate}

\subsection{Energy dispatch and operational KPIs}

Table~\ref{tab:kpis} summarizes operational metrics. The heat pump is the sole thermal supplier across all scenarios (534--537\,GWh annually), with the CHP and boiler not dispatched. Storage provides \SI{163}{GWh} of discharge energy with an average state of charge of 240--\SI{241}{MWh}.

\begin{table}[t]
\centering
\caption{Operational key performance indicators by topology level.}\label{tab:kpis}
\begin{tabular*}{\tblwidth}{@{}lrrr@{}}
\toprule
Metric & \textbf{L1} & \textbf{L2} & \textbf{L3} \\
\midrule
HP heat output [GWh] & 534.4 & 536.5 & 536.9 \\
HP full-load hours [h] & 5,344 & 5,365 & 5,369 \\
HP average COP [--] & 4.00 & 4.00 & 4.00 \\
Storage discharge [GWh] & 162.8 & 163.0 & 163.0 \\
Storage avg.\ SOC [MWh] & 239.9 & 241.2 & 241.4 \\
Grid import [GWh] & 133.5 & 134.0 & 134.1 \\
Network heat loss [GWh] & 0.0 & 2.09 & 2.44 \\
Total CO\textsubscript{2} [kt] & 53.1 & 53.3 & 53.3 \\
\bottomrule
\end{tabular*}
\end{table}

\textbf{Network losses.} L2 predicts \SI{2.09}{GWh} annual pipe losses; L3 predicts \SI{2.44}{GWh} (+17\% relative to L2). The difference arises from L3's more detailed pipe routing (28 pipes vs.\ 4), which captures shorter but more numerous branch pipes with higher surface-area-to-volume ratios.

\textbf{Heat pump dispatch.} The heat pump increases output by \SI{2.5}{GWh} from L1 to L3 to compensate for network losses, corresponding to \SI{0.6}{GWh} additional grid electricity import.

\subsection{Network characteristics}

Table~\ref{tab:network} summarizes the network properties across topology levels. L3's realistic topology with 28 pipes and \SI{14250}{m} total length captures a wider range of pipe diameters (150--\SI{600}{mm}) compared to L2's four trunk lines (350--\SI{500}{mm}).

\begin{table}[t]
\centering
\caption{Network characteristics by topology level.}\label{tab:network}
\begin{tabular*}{\tblwidth}{@{}lrrr@{}}
\toprule
Property & \textbf{L1} & \textbf{L2} & \textbf{L3} \\
\midrule
Nodes & 1 & 5 & 29 \\
Consumer nodes & 0 & 3 & 23 \\
Junction nodes & 0 & 1 & 5 \\
Pipes & 0 & 4 & 28 \\
Total pipe length [m] & 0 & 6,300 & 14,250 \\
Pipe diameter range [mm] & --- & 350--500 & 150--600 \\
Annual pipe loss [GWh] & 0 & 2.09 & 2.44 \\
\bottomrule
\end{tabular*}
\end{table}

\subsection{Computational performance}

Table~\ref{tab:solver} reports solver behavior. L1 solves in approximately 2--3 minutes with an MIP gap below 0.1\%. L2 requires 8--10 minutes (gap $\sim$0.3\%), and L3 solves in 14--20 minutes (gap $<$1\%). The 6$\times$ slowdown from L1 to L3 is acceptable for planning studies.

\begin{table}[t]
\centering
\caption{Computational performance by topology level (HiGHS solver).}\label{tab:solver}
\begin{tabular*}{\tblwidth}{@{}lrrr@{}}
\toprule
Metric & \textbf{L1} & \textbf{L2} & \textbf{L3} \\
\midrule
Continuous variables (K) & $\sim$44 & $\sim$52 & $\sim$56 \\
Binary variables (K) & $\sim$26 & $\sim$28 & $\sim$30 \\
Constraints (K) & $\sim$48 & $\sim$53 & $\sim$61 \\
Solve time [min] & 2--3 & 8--10 & 14--20 \\
MIP gap [\%] & $<$0.1 & $\sim$0.3 & $<$1.0 \\
\bottomrule
\end{tabular*}
\end{table}

\subsection{Dispatch visualization}

Fig.~\ref{fig:dispatch} shows heat dispatch patterns across the three topology levels for a representative winter week. Fig.~\ref{fig:cost} compares the annual cost breakdown. Fig.~\ref{fig:losses} illustrates pipe-level heat losses for L2 and L3. Fig.~\ref{fig:storage} presents storage state-of-charge trajectories.

\begin{figure}[t]
  \centering
  \includegraphics[width=\columnwidth]{fig2_dispatch_comparison.pdf}
  \caption{Comparison of heat dispatch patterns across L1, L2, and L3 for a representative winter week (January 2023). Stacked area chart showing contribution of heat pump and storage to meet demand. Differences between scenarios are minimal, reflecting homogeneous demand distribution.}\label{fig:dispatch}
\end{figure}

\begin{figure}[t]
  \centering
  \includegraphics[width=\columnwidth]{fig3_cost_comparison.pdf}
  \caption{Annual operational cost comparison across topology levels. L1--L3 cost difference (+0.5\%) is driven by grid electricity for loss compensation. Gas fuel cost is zero across all scenarios.}\label{fig:cost}
\end{figure}

\begin{figure}[t]
  \centering
  \includegraphics[width=\columnwidth]{fig4_pipe_losses.pdf}
  \caption{Annual pipe heat losses by topology level. L2 (4 trunk pipes, \SI{6300}{m}) predicts \SI{2.09}{GWh}; L3 (28 pipes, \SI{14250}{m}) predicts \SI{2.44}{GWh}. The 17\% difference arises from L3's finer spatial resolution capturing branch-pipe losses.}\label{fig:losses}
\end{figure}

\begin{figure}[t]
  \centering
  \includegraphics[width=\columnwidth]{fig8_storage_soc.pdf}
  \caption{Thermal storage state-of-charge trajectory over the full year for L1, L2, and L3. Storage behavior is nearly identical across scenarios, with average SOC of 240--\SI{241}{MWh} and peak utilization during winter months.}\label{fig:storage}
\end{figure}


%% ====================================================================
\section{Discussion}\label{sec:discussion}
%% ====================================================================

\subsection{Network topology as a design variable}

The central finding is that for this industrial network with homogeneous demand, network topology abstraction produces modest cost differences (+0.5\% L1$\to$L3). This suggests:

\begin{enumerate}
\item \textbf{The loss model matters more than the topology.} Switching from L1's zero-loss copperplate to L2/L3's physics-based model accounts for approximately 86\% of the L1$\to$L3 cost difference. The additional spatial resolution from L2 to L3 contributes only 14\%.

\item \textbf{L2 captures most of L3's value.} The simplified 5-node model provides $\geq$85\% of L3's operational cost correction at approximately 40\% shorter solve time. This is the key practical guidance for planning engineers.

\item \textbf{Topology matters more for heterogeneous networks.} The small L2$\to$L3 gap reflects a broadly distributed, similar-magnitude load pattern. Networks with remote peaking loads, single-direction corridors, or sheddable industrial loads would show larger topology sensitivity.
\end{enumerate}

\subsection{Fully electrified dispatch}

An unexpected but important finding is that the optimizer selects zero gas boiler dispatch across all scenarios. At a mean COP of 4.0, the heat pump converts grid electricity at an effective cost of $p_\mathrm{el}/\mathrm{COP} \approx \text{\EUR{17.5}/MWh}_\mathrm{th}$ (at \EUR{70}/MWh electricity), which is cheaper than gas at $p_\mathrm{gas}/\eta_\mathrm{boiler} \approx \text{\EUR{50}/MWh}_\mathrm{th}$. This result depends strongly on the COP assumption and electricity/gas price ratio, and the sensitivity to these parameters should be investigated in future work.

\subsection{Unified loss model across abstraction levels}

A key methodological choice is that L2 and L3 use the identical PWL loss formula, varying only in spatial discretization. This isolates topology as the single differentiator: any cost difference is topology-driven, not loss-model-driven. This approach enables fair comparison and clear attribution of cost differences.

\subsection{Practical implications}

Based on the results, we offer the following guidance:
\begin{itemize}
\item \textbf{Screening/feasibility:} L1 copperplate is acceptable if pipe losses are added as an external estimate ($\sim$0.5\% of annual demand for well-insulated networks).
\item \textbf{Detailed planning:} L2 simplified multi-node offers near-L3 accuracy at substantially lower solve time.
\item \textbf{Engineering validation:} L3 realistic topology is justified for final pipe sizing and validation against measured system data.
\end{itemize}

\subsection{Limitations}

\begin{enumerate}
\item \textbf{Homogeneous demand.} Stadtbach serves primarily industrial facilities with similar load profiles. Heterogeneous networks (residential + industrial) would show larger topology effects.
\item \textbf{Single year.} The 2023 dataset represents one year. Multi-year studies are recommended for robust conclusions.
\item \textbf{Fixed COP method.} All scenarios use identical analytical LMTD. Future work should validate against real heat pump performance data.
\item \textbf{No demand-side management.} Flexible loads are not modeled; such assets would increase topology sensitivity.
\item \textbf{No pressure-flow coupling.} The Darcy--Weisbach nonlinearity ($\Delta p \propto v^2$) is handled post-processing, not within the MILP.
\end{enumerate}


%% ====================================================================
\section{Conclusion}\label{sec:conclusion}
%% ====================================================================

This study addresses a practical but under-quantified question in district heating system planning: \emph{How much does the level of network topology detail affect operational dispatch optimization outcomes, when the loss physics are held identical across all models?}

We compared three topology abstraction levels---L1 copperplate (no spatial structure), L2 simplified 5-node (zonal aggregation), and L3 detailed 29-node (realistic network)---applied to the Stadtbach industrial heat network in Austria (\SI{76}{MW} peak, \SI{517}{GWh/year}). By fixing all asset capacities and using identical COP calculation and pipe-loss physics (PWL approximation of $Q = U \cdot L \cdot \Delta T$) at L2 and L3, network topology was the sole experimental variable.

Main findings:

\begin{enumerate}
\item \textbf{Including physical losses matters; topology detail less so.} The L1$\to$L3 annual cost increase is 0.5\% (\EUR{74,195}). Of this, 86\% is attributable to including physics-based losses (L1$\to$L2), while only 14\% comes from additional spatial resolution (L2$\to$L3).

\item \textbf{L2 is sufficient for planning.} The 5-node simplified model captures $>$85\% of L3's cost correction at approximately 40\% shorter solve time---the key practical result.

\item \textbf{Fully electrified dispatch emerges.} At COP~$\approx$~4.0, the heat pump dominates all scenarios with zero gas boiler dispatch. This finding is COP- and price-dependent.

\item \textbf{Computational efficiency is practical.} L3 MILP (29 nodes, 8,760 hours) solves in under 20 minutes, enabling full-year topology studies within a half-day of computation.

\item \textbf{Framework availability.} The CALION framework provides all three topology levels via configuration files, with built-in sensitivity analysis and benchmarking modules.
\end{enumerate}

Future research should: (i)~extend to heterogeneous networks to characterize when topology sensitivity increases; (ii)~conduct multi-year studies; (iii)~incorporate investment optimization to assess capacity sizing changes with topology level; and (iv)~develop pressure-constrained MILP extensions.


%% ====================================================================
%% Data availability
%% ====================================================================
\section*{Data availability}

The Stadtbach metering data are subject to a confidentiality agreement with the network operator and cannot be made publicly available. The CALION source code, configuration files for all three topology scenarios (L1/L2/L3), and a synthetic benchmark dataset with equivalent statistical properties are publicly available at \url{https://github.com/[repository-to-be-added]}.


%% ====================================================================
%% Declaration
%% ====================================================================
\section*{Declaration of competing interest}

The authors declare that they have no known competing financial interests or personal relationships that could have appeared to influence the work reported in this paper.


%% ====================================================================
%% Acknowledgements
%% ====================================================================
\section*{Acknowledgements}

%% TODO: Add funding sources, grant numbers, and acknowledgements.
%% Example: This research was funded by the Austrian Research Promotion Agency (FFG) under grant [number].


%% ====================================================================
%% CRediT
%% ====================================================================
\printcredits


%% ====================================================================
%% Appendix
%% ====================================================================
\appendix

\section{Linearization proofs}\label{app:proofs}

\begin{theorem}[COP pre-computation preserves MILP]\label{thm:cop}
Let $\mathrm{COP}^*[t]$ be computed offline by interpolating a 2D table at exogenous temperatures. Then $Q_\mathrm{hp}[t] = \mathrm{COP}^*[t] \cdot P_\mathrm{hp}[t]$ is linear in $P_\mathrm{hp}[t]$ and the MILP remains tractable.
\end{theorem}

\noindent\textit{Proof.} Since $T_\mathrm{source}[t]$ is exogenous, $\mathrm{COP}^*[t]$ is a fixed scalar for each $t$. The constraint $Q_\mathrm{hp}[t] - \mathrm{COP}^*[t] \cdot P_\mathrm{hp}[t] = 0$ is a linear equality in $(Q_\mathrm{hp}[t], P_\mathrm{hp}[t])$. The resulting MILP has $\sim$56,000 continuous and $\sim$30,000 binary variables with sparsity $>$99.97\%, efficiently solvable by branch-and-cut~\citep{Wolsey1998}. \hfill$\square$

\textbf{Error bound.} A 4\% COP error (combined interpolation + measurement) propagates to $\leq$1.4\% total system cost error for heat-pump-dominated systems.

\begin{theorem}[PWL approximation error bound]\label{thm:pwl}
For a $C^2$ loss function $f: [a,b] \to \mathbb{R}$ with $N$ uniformly-spaced PWL breakpoints:
\begin{equation}
\|f - f_\mathrm{PWL}\|_\infty \leq \frac{M(b-a)^2}{8N^2}
\end{equation}
where $M = \max_{x} |f''(x)|$.
\end{theorem}

\noindent\textit{Application.} For a \SI{500}{MWh} stratified tank ($M \approx 0.1$\,W/K$^2$, $N=10$): error $\leq \SI{31}{W} < 1$\% of typical losses. \hfill$\square$

\begin{theorem}[Big-M constraint tightness]\label{thm:bigm}
The mutual exclusivity constraints (Eqs.~\ref{eq:grid_buy}--\ref{eq:grid_sell}, \ref{eq:stor_charge}--\ref{eq:stor_discharge}) with $M \geq \max(P_\mathrm{buy,max}, P_\mathrm{sell,max})$ achieve zero relaxation gap at integer optimum.
\end{theorem}

\noindent\textit{Proof.} The constraints force exactly one of $\{P_\mathrm{buy}, P_\mathrm{sell}\}$ (or $\{Q_c, Q_d\}$) to zero. No fractional binary solution exists at optimum. \hfill$\square$


%% ====================================================================
%% References
%% ====================================================================
\bibliographystyle{cas-model2-names}
\bibliography{references}

\end{document}
